\begin{textsample}{POS Dim 3 – pa – Score 15.00 – pa/pa\_inf\_4.txt}  \label{ex:f3_pos_235}
3.1.2.1 Concepções de \textbf{Infâncias} O conceito de \textbf{infância}, enquanto uma categoria social e fase específica da existência humana é parte de uma discussão muito recente em termos historiográficos.
Ariès ( 1981 ) foi um dos precursores no desenvolvimento de estudos acerca da noção moderna de \textbf{infância}, contribuindo para o iniciar de discussões que convergem para o entendimento de que o sujeito \textbf{criança} vem sendo visto de diferentes maneiras ao longo dos tempos até chegar à compreensão que se tem como referência atualmente.
Para Heywood ( 2004 ), a preocupação pelo período da \textbf{infância} é um fenômeno recente, muito por conta dos poucos \textbf{registros} que se têm sobre as memórias e experiências de \textbf{infância} em épocas passadas, pois havia limitado interesse em conhecer as especificidades dos sujeitos nessa etapa da vida.
Na sociedade medieval, por exemplo, a \textbf{centralidade} em assuntos religiosos retirou muitos temas do eixo de interesse da época e a \textbf{infância} foi um deles, predominando o foco na vida \textbf{adulta}.
A modernidade trouxe então a ideia de que a \textbf{criança}, tida como ingênua, pura e ociosa, fazia parte de um período que precisava ser alvo de investimentos morais, educacionais e de \textbf{cuidados} com a saúde.
Com isso se promoveu um modelo de \textbf{infância} universal que foi divulgado e projetado a partir do padrão burguês de \textbf{criança}, com base em critérios de idade e dependência do \textbf{adulto}, característicos de um tipo específico de papel social por ela assumido no interior dessa classe.
O olhar atual sobre a \textbf{infância} é uma construção forjada na modernidade.
Segundo Stearns ( 2006 ), nesse contexto, a \textbf{infância} vai englobar três questões essenciais e inter-relacionadas, que irão influenciar e promover um novo modo de ver as \textbf{crianças}.
A primeira envolve a passagem da \textbf{infância}, até então voltada ao trabalho, para a escolaridade ; a segunda diz respeito à decisão de limitar, levando em consideração os altos custos, o tamanho das famílias a patamares mais baixos ; e, por fim, a redução da taxa de mortalidade \textbf{infantil}.
Na atualidade, muitos estudos são realizados no sentido de pensar a \textbf{infância} como uma categoria heterogênea que vê as \textbf{crianças} como sujeitos sociais e historicamente situadas em determinado contexto, e constituídos pelas interações e experiências vivenciadas nas suas realidades, o que implica dizer que seu desenvolvimento se dá entre outros seres humanos, em um tempo e espaço determinado.
A \textbf{infância}, portanto, é um produto que se constitui a partir de um conjunto de características que possibilita pensar que ela não existe somente de uma forma nem vivencia as mesmas realidades, haja vista que nem toda \textbf{infância} é repleta de tempo livre para \textbf{brincar}, de ausência de responsabilidades \textbf{adultas} e de direitos à saúde e educação de qualidade assegurada. [... ] sujeito individual que carrega desde o \textbf{nascimento} as expectativas sociais e ao desvendar o mundo e mergulhado nele aprende ou pode aprender a se constituir indivíduo, alegoricamente como espécie de um cristal, é pedra, pois sedimentada pela formação ; todavia o desenho toma forma própria ( SOUZA, 2007, p. 74 ).
Nessa perspectiva, a \textbf{infância} não pode ser entendida de maneira homogênea, afinal não existe uma maneira exclusiva de vivê -la, posto que há apenas uma única, mas várias \textbf{infâncias} ( FREITAS ; KUHLMAM JR, 2002 ), que devem ser visualizadas a partir de suas especificidades econômicas, sociais e culturais ; é possível, desse modo, falar em \textbf{infância} pobre, rica, oriental, ocidental, urbana, agrária, indígena, ribeirinha, quilombola, etc.
Cada qual constituída por características que se aproximam e se distanciam entre si e seus próprios \textbf{pares}.
A \textbf{infância} precisa ser entendida também como fase da vida em que os sujeitos que nela se encontram não sejam tomados como um projeto a ser concretizado, o vir a ser, o que ainda não é e precisa ser preenchido para deixar de ser incompleto.
Essa visão foi, historicamente, alicerçada no próprio significado etimológico do termo \textbf{infância} que, conforme Lajolo ( 2011 ), tem origem na língua latina : infante ( in : prefixo que indica negação ; fante : particípio presente do verbo latino fari, que significa falar, denotando a ideia de ausência de fala ).
Partindo desse princípio, a \textbf{criança} por muito tempo foi concebida como um ser que não fala, logo, que não tem como produzir a partir do seu viés sua própria história, e [... ] por não falar, a \textbf{infância} não se fala e, não se falando, não ocupa a primeira pessoa nos discursos que dela se ocupam.
E, por não ocupar esta primeira pessoa, isto é, por não dizer eu, por jamais assumir o lugar de sujeito do discurso, e, consequentemente, por consistir sempre um ele/ela nos discursos alheios, a \textbf{infância} é sempre definida de fora ( LAJOLO, 2011, p. 230 ).
O silenciamento da \textbf{criança} sinaliza sobre ela um posto de subalternidade, aqui compreendida quando há a supremacia de um sujeito em detrimento de outro na medida em que lhe é negada as instâncias de fala.
Spivak ( 2010, p.126 ) afirma que tal negação, resultante de processos hegemônicos, orquestram a negação de representação, de dialogismo e de uma participação ativa e política do sujeito na sociedade.
A \textbf{criança}, portanto, se compreendida de maneira inferior ao \textbf{adulto}, se calada, já que ‘ o subalterno não pode falar ’, tem negada a sua voz, por meio da qual manifesta suas formas de pensar e, consequentemente, seu ‘ existir ’ social.
Para Freitas ( 2007, p. 90 ), as \textbf{crianças} parecem basicamente ser aquilo o que delas se fala, uma vez que “ são os incapazes em relação aos capazes ; são os ociosos em relação aos produtivos, são os normais em relação aos anormais ”.
Entretanto, essa visão limita o sujeito \textbf{criança}, fazendo desse alheio e estático perante aos fatos e condicionantes que emergem em seus contextos reais de interação.
Esse pensamento perdurou por muito tempo inclusive em pesquisas científicas : a \textbf{criança} era tratada como objeto a ser medido, observado, descrito, analisado e interpretado, ou seja, sempre como “ o outro ” em relação àquele que a nomeia e a estuda.
Se a \textbf{criança} é aquela que não fala, está em desenvolvimento, é incompleta e não tem o que falar, não apresenta capacidade de expressar suas particularidades, ela apenas imita, reproduz o \textbf{adulto}.
Martins Filho ( 2006 ) verifica na referência do estado, ao investigar os processos de socialização entre \textbf{crianças} e entre elas e os \textbf{adultos}, constatou que a forma como os \textbf{adultos} percebem as \textbf{crianças} reflete nas relações e nos modos como estes se dirigem a elas.
Para o autor, se o \textbf{adulto} considera a \textbf{criança} como ator social ele a \textbf{ouve} e respeita as suas especificidades e suas manifestações culturais, mas se o \textbf{adulto} tem uma visão de \textbf{criança} como sujeito padronizado, continuará tratando -a como ser sem \textbf{vontades} próprias, incompleto, moldável e apenas imitador de práticas culturais alheias.
Porém hoje a \textbf{criança} é vista não somente como sujeito de direitos, a partir de toda a legislação oficial que ampara e legitima suas necessidades, mas também como produtora de cultura e não mais como simples reprodutora das manifestações realizadas no universo \textbf{adulto}, sendo assim ela é concebida como Sujeito histórico e de direitos que, nas interações, relações e práticas cotidianas que vivencia, constrói sua identidade pessoal e coletiva, \textbf{brinca}, imagina, fantasia, deseja, aprende, observa, experimenta, narra, questiona e constrói sentidos sobre a natureza e a sociedade, produzindo cultura ( BRASIL, 2010a, p. 12 ).
O sujeito \textbf{criança}, ao ser tratado como protagonista no seu processo de socialização e interação com o mundo, a partir das relações entre seus \textbf{pares} e com os \textbf{adultos}, constrói interpretações particulares, apresentando certa autonomia para estabelecer significados de suas vivências.
Como bem afirma Sarmento ( 2004 ), as \textbf{crianças} experimentam a cultura em que se inserem distintamente da cultura \textbf{adulta}, produzindo uma que lhes é própria, logo, a questão fundamental no estudo das culturas \textbf{infantis} é a interpretação da sua autonomia em relação aos \textbf{adultos}.
É justamente nesse contexto, em privilegiar a \textbf{escuta} da \textbf{criança} e compreender as suas culturas, que várias áreas do conhecimento, como a Psicologia, a Sociologia, a Antropologia e a Educação vêm se direcionando para a \textbf{infância}, uma vez que já se reconhece que desde a mais tenra idade, nas suas interações sociais, os sujeitos vão somando “ impressões, gostos, antipatias, \textbf{desejos}, medos etc., desenvolvendo sentimentos e \textbf{percepções}, cada vez mais diversificados e definidos, atribuindo significados, construindo a sua identidade ” ( CRUZ, 2008, p. 13 ).
Compreender e dar visibilidade as \textbf{infâncias} representa o iniciar de sua valorização e reconhecimento enquanto categoria social, pois, as \textbf{crianças} têm muito a dizer sobre as suas formas de ver o mundo, sobre preconceitos, sobre o poder e a autoridade que os \textbf{adultos} exercem sobre elas ( QUINTEIRO, 2002 ).
É necessário conhecer mais sobre as culturas \textbf{infantis}, os modos de vida das \textbf{crianças}, as \textbf{crianças} que frequentam a escola, como aprendem, o que aprendem, o que sentem e o que pensam.
Nesse sentido, ao se tomar as múltiplas \textbf{infâncias} vividas em contextos heterogêneos, o entendimento de que os sujeitos que delas fazem parte são \textbf{crianças} concretas, vivas, reais e que tem algo a dizer a partir de seus olhares, contribui, sobremaneira, para subsidiar ações de outras pessoas e entidades que possibilitam a elaboração de currículos, práticas e programas que tomam como ponto de partida as \textbf{crianças} e suas especificidades, para que de fato elas tenham condições de usufruir de suas \textbf{infâncias}.

% matched lemmas: adulto, brincar, centralidade, criança, cuidado, desejo, escuta, infantil, infância, nascimento, ouvir, par, percepção, registro, vontade
\end{textsample}
